\documentclass[12pt,a4paper]{article}
\usepackage[utf8]{inputenc}
\usepackage[T1]{fontenc}
\usepackage[polish]{babel}
\usepackage{graphicx}
\usepackage{geometry}
\geometry{a4paper, margin=1.5cm}
\usepackage{hyperref}
\usepackage{float}
\usepackage{listings}
\usepackage{xcolor}

\begin{document}

\begin{titlepage}
    \centering
    \vspace*{2cm}
    {\Large \textbf{Laboratorium Architektury Systemów Komputerowych}} \\[1.5cm]
    {\Huge \textbf{Sprawozdanie z projektu nr 2}} \\[1cm]
    {\Large Aplikacja do pomiaru czasu reakcji} \\[3cm]
    \begin{flushleft}
        \textbf{Wykonał:} Łukasz Mroczek 198146, Łukasz Orluk 197641\\       
        \textbf{Grupa:} KSD2A\\
    \end{flushleft}
    \vfill
    {\large \today}
\end{titlepage}

\newpage
\tableofcontents

\newpage

\section{Sformułowanie zadania i założenia szczegółowe}

Celem zadania było stworzenie aplikacji uzależnionej od czasu, umożliwiającej pomiar czasu reakcji użytkownika na różnego rodzaju bodźce. Program miał symulować tester sprawności psychomotorycznej, np. dla kandydatów na kierowców.

\subsection{Założenia szczegółowe}

\begin{itemize}
    \item Aplikacja została wykonana w języku Python z wykorzystaniem biblioteki PyQt6 oraz matplotlib.
    \item Program bada trzy typy reakcji:
    \begin{itemize}
        \item prosty czas reakcji na bodziec optyczny,
        \item prosty czas reakcji na bodziec akustyczny,
        \item złożony czas reakcji (wybór odpowiedniego klawisza).
    \end{itemize}
    \item Każdy test składa się z:
    \begin{itemize}
        \item opisu,
        \item fazy treningowej (bez oceny),
        \item testu właściwego (z pomiarem wyników).
    \end{itemize}
    \item Czas reakcji mierzony jest przy użyciu funkcji \texttt{time.perf\_counter()}, zapewniającej wysoką dokładność.
    \item Pomiędzy próbami stosowane są losowe opóźnienia (1–4 s), aby zapobiec przewidywaniu momentu pojawienia się bodźca.
    \item Wyniki prezentowane są w formie:
    \begin{itemize}
        \item wartości liczbowych (średnia, min, max, mediana),
        \item wykresu słupkowego.
    \end{itemize}
\end{itemize}

\section{Opis przyjętych rozwiązań programowych}

\subsection{Struktura aplikacji}

Aplikacja została podzielona na kilka głównych komponentów:

\begin{itemize}
    \item \textbf{MainWindow} – zarządza przełączaniem ekranów,
    \item \textbf{IntroScreen} – wyświetla instrukcje,
    \item \textbf{TestRunnerScreen} – realizuje test,
    \item \textbf{ResultsScreen} – prezentuje wyniki,
    \item \textbf{models.py} – przechowuje dane i statystyki.
\end{itemize}

---

\subsection{Obsługa czasu i bodźców}

Do sterowania testem wykorzystano klasę \texttt{QTimer}. Timer generuje opóźnienie przed pojawieniem się bodźca:

\begin{lstlisting}[language=Python, caption=Losowe opóźnienie przed bodźcem]
def random_delay_ms():
    return random.randint(1500, 4000)
\end{lstlisting}

Po upływie czasu wywoływana jest funkcja generująca bodziec.

---

\subsection{Pomiar czasu reakcji}

Czas reakcji mierzony jest w następujący sposób:

\begin{lstlisting}[language=Python, caption=Pomiar czasu reakcji]
def reaction_time_ms(start_time):
    return (time.perf_counter() - start_time) * 1000
\end{lstlisting}

Zapamiętywany jest moment pojawienia się bodźca, a następnie odejmowany od czasu reakcji użytkownika.

---

\subsection{Obsługa klawiatury}

Program umożliwia reakcję zarówno przyciskiem, jak i klawiaturą:

\begin{lstlisting}[language=Python, caption=Obsługa klawiatury]
def keyPressEvent(self, event):
    if event.key() == Qt.Key.Key_Space:
        self.on_space_response()
\end{lstlisting}

W testach złożonych wykorzystywane są klawisze L oraz P.

---

\subsection{Zbieranie wyników}

Każda próba zapisywana jest jako obiekt klasy \texttt{TrialResult}. Wyniki całego testu przechowywane są w klasie \texttt{TestSummary}.

\begin{lstlisting}[language=Python, caption=Filtracja poprawnych wyników]
def valid_times(self):
    return [
        result.reaction_time_ms
        for result in self.real_results
        if result.correct and result.reaction_time_ms is not None
    ]
\end{lstlisting}

---

\subsection{Obliczanie statystyk}

Na podstawie poprawnych wyników obliczane są statystyki:

\begin{lstlisting}[language=Python, caption=Obliczanie statystyk]
def stats(self):
    return {
        "avg": statistics.mean(times),
        "min": min(times),
        "max": max(times),
    }
\end{lstlisting}

Dodatkowo liczona jest mediana oraz odchylenie standardowe.

---

\section{Dyskusja osiągniętych wyników}

\subsection{Zalety aplikacji}

\begin{itemize}
    \item Dokładny pomiar czasu dzięki funkcji \texttt{perf\_counter}.
    \item Zastosowanie losowych opóźnień zwiększa wiarygodność testu.
    \item Intuicyjny interfejs użytkownika.
    \item Obsługa zarówno myszy, jak i klawiatury.
    \item Graficzna prezentacja wyników (wykres).
\end{itemize}

\subsection{Wady aplikacji}

\begin{itemize}
    \item Brak zapisu wyników do pliku.
    \item Brak możliwości powtórzenia pojedynczego testu bez restartu aplikacji.
    \item Ograniczona liczba typów bodźców.
\end{itemize}

\section{Podsumowanie}

Zrealizowana aplikacja spełnia wymagania zadania. Program poprawnie mierzy czas reakcji użytkownika oraz prezentuje wyniki w sposób czytelny i przejrzysty. Zastosowane rozwiązania programowe są proste, ale skuteczne i mogą stanowić podstawę do dalszej rozbudowy aplikacji.

\end{document}